\scanpage{54}
1) $\mathcal M_{\xi_0}$ a \inserted{la structure d'} un groupe formel
lisse sur $\deleted{R}\operatorname{Spf}\Lambda$,
grâce à la structure de groupe du foncteur qu'il représente.
Le groupe formel en question est un \uline{tore formel},
\deleted{\unsure{il est déterminé}} \inserted{donc déterminé}
par sa réduction sur $k$, et cette dernière est le tore
\inserted{formel tordu} $T_p(M_0)^\vee\otimes N_0$.

2) $\mathcal M_\xi$ est un \deleted{fibré principal homogène}
\inserted{torseur} sous $\mathcal M_{\xi_0}$.
\inserted{[\unsure{qui est en substance une première version}
\unsure{provisoire|corollaire} du th. de Tate affirmant qu'on peut tjrs lifter
\unsure{l'inf\textsuperscript{t}|l'inf.} d'un groupe \deleted{formel} $p$-divisible]}

Ne supposons plus $N_0$ de t.m., mais \unsure{seult radiciel}.
On veut préciser la structure de la variété formelle des
modules pour $E_0$. [NB. $E_0$ détermine \unsure{uniqt}
$M_0$, $N_0$ et la structure d'extension!)

On a un hom naturel
\[
X=\mathcal M_{\unsure{E_0}}\longrightarrow\mathcal M_{N_0}=Y
\]
Ce dernier fait de $\mathcal M_{E_0}$ un
\deleted{fibré} torseur \deleted{principal homogène}
sous $\mathcal N_Y\otimes_{\mathbb Z_p}T_p(E_0)^\vee_Y$.
Donc la structure de $\deleted{X}\mathcal M_{E_0}$ est ramenée
à celle de $\mathcal M_{N_0}$.
[Dans le cas $N_0$ de t.m., $Y$ est évident: un point].
\margin{64}
